% ========== 1.1 PROBLEM DEFINITION ==========

\section{Problem Definition}
\label{sec:problem-definition}

\lettrine{M}odern development environments face significant challenges when integrating artificial intelligence capabilities. Traditional IDEs were designed decades ago with the assumption that human developers would be the primary actors in all development activities. This human-centric design creates fundamental barriers when attempting to incorporate AI features that require deep system integration.

Current IDEs suffer from performance issues when AI capabilities are added. Extension latency increases substantially, memory consumption grows dramatically, and startup times extend significantly. These problems stem from architectural decisions made before AI integration was a consideration. The systems were simply not designed to handle the asynchronous nature of AI operations or the variable resource requirements of intelligent agents.

The retrofit approach of adding AI features to existing architectures creates additional problems. Developers experience fragmented workflows where AI tools operate in isolation without sharing context. Context switching between traditional coding interfaces and AI assistance breaks cognitive flow. Manual coordination between different tools creates friction that reduces productivity rather than enhancing it.

The core issue is that current development environments treat AI as supplementary tooling rather than foundational architecture. This mismatch prevents effective human-AI collaboration and limits what AI can contribute to the development process. A new approach is needed that designs development environments specifically for AI integration from the ground up.
